\section{An efficient parameterized approximation scheme for treewidth}\label{sec:tw}

In this section we describe the main result of this paper which is an EPAS for
\SF parameterized by treewidth. We begin by giving two preliminary tools
(\autoref{lem:nice-tree-decomp} and \autoref{lem:aspect-ratio}) which
facilitate the algorithm by ensuring that the given tree decomposition has
logarithmic height and that the instance has aspect ratio (ratio of the weights
of the heaviest over the lightest edge) bounded by a polynomial in $n$.

We then go on to \autoref{sec:tw-height} where we introduce a second parameter,
the height $h$ of the decomposition. Our goal is to fix an almost-optimal
solution $F_\eps$ and describe an algorithm that produces a partition $\zeta_B$ of the active terminals
for each bag $B$ of the decomposition, where $\zeta_B$ is
\emph{a refinement} of the partition implied by $F_\eps$
(\autoref{lem:near-opt-1}). In other words, we seek a partition $\zeta_B$ of
$A_B$ such that if two terminals $t_1,t_2$ are in the same set of $\zeta_B$,
then they are also in the same component of $F_\eps$. Of course, it is trivial
to achieve this by giving a $\zeta_B$ where each active terminal is in its own
set, so the interesting part here is how we group terminals together in a way
that in the end allows us to bound $|\zeta_B|$ by a polynomial of
$k+h+\frac{1}{\eps}+\log n$, while still ensuring that $F_\eps$ is almost
optimal. 

The partition $\zeta_B$ of \autoref{lem:near-opt-1} is not yet conforming,
because two terminals which are in distinct sets of $\zeta_B$ may still be in
the same component of $F_\eps$, and thus we cannot apply \cref{thm:tw-dp} at this point. 
Therefore in \autoref{sec:tw-log-height}, given $\zeta_B$ we
focus on how to obtain every possible partition of the set of
active terminals, which could be conforming with an almost-optimal solution. By
an appropriate use of $\delta$-nets, similar to \cite{bateni2011approximation},
we are able to ``guess'' (that is, brute-force) a choice of a small number of
net points per active component. Since the number of choices for each point is
at most $|\zeta_B|$ and we choose roughly $O(k^2/\epsilon)$ points in total, the
total number of produced partitions (and hence the running time given by
\autoref{thm:tw-dp}) is of the form $(\log n + k
+\frac{1}{\epsilon})^{O(k^2/\epsilon)}n^{O(1)}$, which is FPT.

Let us now recall a result of \citet{BodlaenderH98} which states that a tree
decomposition of logarithmic height can always be obtained.

\begin{lemma}[\cite{BodlaenderH98}]\label{lem:nice-tree-decomp}
%
%
Given a tree decomposition of width $k$ of a graph $G$ on $n$ vertices, there is a
polynomial time algorithm computing a nice tree decomposition of $G$ of
width $O(k$) and height $O(k\log n)$.  \end{lemma}

%
\begin{proof}%
It is shown in \cite{BodlaenderH98} that any tree decomposition
of width $k$ can be transformed in polynomial time into a tree decomposition of
width $O(k)$ and height $O(\log n)$ where the tree of the decomposition has
maximum degree $3$. It is now not hard to make this decomposition nice by
replacing all nodes with two children by Join nodes, and by inserting between
any node and its parent a sequence of $O(k)$ new nodes so that the symmetric
difference between any node and its parent contains at most one vertex.
\end{proof} 
%

We also need to reduce the aspect ratio of the given graph to a polynomial.
This can be done using a standard technique, where however we need to make
sure that the treewidth of the given graph remains bounded.  
Note that the aspect ratio of the resulting graph $G'$ in the following lemma
is polynomially bounded in the size of the original graph, 
but not necessarily in the size of $G'$  (because $G'$ may have
significantly fewer vertices).
%
%
%
%
%
%

\begin{lemma}\label{lem:aspect-ratio} %
Given $\eps>0$, an instance of \SF on a
graph $G$ with $n$ vertices, and a (nice) tree decomposition $T$ of width $k$
and height $h$ for $G$, in polynomial time we can compute an instance on a
graph $G'$ with at most $n$ vertices and a (nice) tree decomposition~$T'$ of
width at most $k$ and height $h$ for~$G'$, such that the ratio of the longest
to the shortest edge in~$G'$ is at most $2n/\eps$, and any
$\alpha$-approximation for $G'$ can be converted into an
$(\alpha+\eps)$-approximation for $G$.  
\end{lemma} 

%
\begin{proof}%
The first
step is to compute a $2$-approximation $F_2$ for \SF in $G$, using the
polynomial time algorithm of~\cite{agrawal1991trees}.  The new graph $G'$ is
obtained from $G$ by first removing all edges of length more than $\cost(F_2)$
and then contracting every edge of length less than
$\frac{\eps}{2n}\cost(F_2)$, where $n$ is the number of vertices of $G$.  If a
contracted edge was incident to a terminal, then the new vertex is declared a
terminal and the demands are updated correspondingly (note that this may
introduce trivial demands from the new terminal to itself if a demand pair is
connected by a path of edges being contracted).  We modify the tree
decomposition $T$ to obtain $T'$ as follows: whenever we contract the endpoints
of an edge $v_1v_2$ into a new vertex $w$, we replace all occurrences of $v_1$
and $v_2$ in $T$ by $w$. It is not hard to see that this keeps a valid decomposition
of the same height and can only decrease the width. Furthermore, if the
original decomposition was nice, the new decomposition can easily be made nice, 
if we contract every bag $B$ with a unique child $B'$
whenever $B=B'$ (which were introduce or forget nodes previously). 
Also, clearly the ratio between longest and shortest edge in
$G'$ is at most $2n/\eps$.

It remains to show that an $\alpha$-approximate solution $F'_\alpha$ in $G'$ is not distorted by much when converting it 
from $G'$ to $G$. Starting with $F'_\alpha$ the conversion is simply done by iteratively uncontracting those edges that were contracted to obtain $G'$ from $G$: 
if the solution becomes infeasible after uncontracting some edge $e$ we just add it to the solution to make it feasible again. 
Let $F$ denote the solution obtained for $G$ from $F'_\alpha$, and note that less than $n$ edges are added to $F'_\alpha$ in this process, as $F$ is a forest in~$G$.
This means that $\cost(F)<\cost(F'_\alpha)+\frac{\eps}{2}\cost(F_2)$ since every contracted edge has length less than $\frac{\eps}{2n}\cost(F_2)$.
Now consider an optimum solution $F^\star$ in $G$. It can be converted into a solution of cost at most $\cost(F^\star)$ in $G'$ by contracting all edges of length less than~$\frac{\eps}{2n}\cost(F_2)$, 
since $F^\star$ cannot contain any of the removed edges of length more than $\cost(F_2)\geq\cost(F^\star)$.
Thus the optimum of $G'$ has cost at most $\cost(F^\star)$, and because~$F'_\alpha$ is an $\alpha$-approximation in $G'$ we get $\cost(F'_\alpha)\leq\alpha\cost(F^\star)$. 
At the same time, $\cost(F_2)\leq 2\cost(F^\star)$, which together with the previous inequality gives 
$\cost(F)<\alpha\cost(F^\star)+2\frac{\eps}{2}\cost(F^\star)=(\alpha+\eps)\cost(F^\star)$, which concludes the proof.
\end{proof}    
%

For simplicity, in the following we will scale the edge lengths of any given graph so that the shortest edge has length $1$.
In particular, after applying \autoref{lem:aspect-ratio}, the longest edge has length at most~$2n/\eps$.

\subsection{Tree decompositions with bounded height}\label{sec:tw-height}

In this section we informally assume that the height $h$ of the given tree
decomposition is bounded as well as the width $k$. 
%
Our aim is to prove the following statement, where we restrict ourselves to input graphs of polynomial aspect ratio, which we may do according to \autoref{lem:aspect-ratio} (keeping in mind that $n$ is the number of vertices of the original input graph).

\begin{lemma}\label{lem:near-opt-1} 
Let an instance of \SF on a graph $G$ with at most $n$ vertices 
be given together with a tree decomposition~$T$ of
width~$k$ and height $h$ for $G$. For any $\eps>0$, if the ratio between the
longest and shortest edge of $G$ is at most $2n/\eps$, then there exists a
$(1+\eps)$-approximation $F_\eps$ with the following properties. There exists a
polynomial time algorithm, which for every bag $B$ of~$T$ outputs a partition
$\zeta_B$ of the active terminals~$A_B$, such that each set of~$\zeta_B$
belongs to the same component of $F_\eps$ and
$|\zeta_B|=O(\frac{k^4h^2}{\eps^2}\log\frac{n}{\eps})$.  
\end{lemma}

%
%

To prove \autoref{lem:near-opt-1} we first identify the solution $F_\eps$, after which we will show how to compute the partitions $\zeta_B$.

\subsubsection{A near-optimal solution}

The high-level idea to obtain a $(1+\eps)$-approximate solution $F_\eps$ is to connect components of the optimum solution~$F^\star$ that lie very close to each other. In particular, if the distance between two components $C$ and~$C'$ of $F^\star$ is of the form $f(k,h,\eps)\cost(C)$ for some small enough function $f$, then we may hope to add a shortest path between $C$ and $C'$ and charge this additional cost to $C$, in order to obtain a $(1+\eps)$-approximation. Unfortunately, this approach is not viable, since the number of components that are very close to $C$ may be very large, meaning that the function $f$ in the distance bound would have to linearly depend on the number of vertices in order to result in a $(1+\eps)$-approximation. This in turn would mean that the size of the partition $\zeta_B$ would depend polynomially on the number of vertices, making it unsuitable for an FPT time algorithm. This issue lies at the heart of the problem and is the reason for why it is non-trivial to obtain an approximation scheme parameterized by the treewidth. To get around this issue we will measure the distance between components using a modified cost function, which we define next.

Given a bag $B$ of the rooted tree decomposition $T$, we denote by $T_B$
the subtree of $T$ rooted at the node associated with $B$, and by $G_B=G[V_B]$ the graph induced by 
the vertices $V_B$ lying in bags of $T_B$. We also define the graph 
$G^\downarrow_B\subseteq G_B$ as the graph spanned by all edges of $G_B$, 
except those induced by~$B$, i.e., the edge set of $G^\downarrow_B$ is
\[
E(G^\downarrow_B)=\{uv\in E(G_B)\mid u\notin B\lor v\notin B\}.
\] 
The cost of a component $C$ of some \SF solution restricted to $G^\downarrow_B$
only counts the edge weights of $C$ in $G^\downarrow_B$, and is denoted by
\[
\cost^\downarrow_B(C)=\sum_{e\in E(C)\cap E(G^\downarrow_B)}w(e).
\]

Based on these definitions, we fix an optimal solution $F^\star$ and construct
a solution $F_\eps$ by initially setting $F_\eps=F^\star$, and then connecting
components by exhaustively applying the following rule, where we say that two
components $C$ and $C'$ \emph{share} a bag $B$ if $V(C)\cap B\neq\emptyset$ and
$V(C')\cap B\neq\emptyset$:

\begin{description}
    \item[Rule 1:]\customlabel{rule1}{Rule~1} if $C,C'$ are components of $F^\star$ sharing a bag $B$ with $\dist(C,C')\le\frac{\eps}{kh}\cdot\cost^\downarrow_B(C)$ but $C$ and~$C'$ are in different components of $F_\eps$, then add a shortest path of length $\dist(C,C')$ between $C$ and~$C'$ to the solution~$F_\eps$.
\end{description}

\begin{lemma}\label{lem:feps}
The cost of the solution $F_\eps$ obtained by \autoref{rule1} from $F^\star$ is at most $(1+\eps)\cost(F^\star)$.
\end{lemma}

\begin{proof}
It suffices to prove that the cost of all paths added to $F^\star$ in order to obtain $F_\eps$ according to \autoref{rule1}
is at most $\eps\cdot\cost(F^\star)$.
For this we use a charging scheme that charges new paths to components of $F^\star$. 
In particular, we charge a path of length $\dist(C,C')\le\frac{\eps}{kh}\cdot\cost^\downarrow_B(C)$ to component~$C$.

Fix a component $C$ of $F^\star$ and a bag $B$ with $V(C)\cap B \neq \emptyset$. We define $\textrm{charge}(C,B)$ to be the cost we charge to $C$ for operations involving other components of $F^\star$ that share $B$. It is not hard to see that $\textrm{charge}(C,B)\le \frac{\eps}{h}\cdot \cost^\downarrow_B(C)$, because there are at most $k$ other components of $F^\star$ that share $B$. 

For $\ell\in\{0,\ldots,h-1\}$, let $\mathcal{B}_\ell$ be the set of bags of the tree decomposition that appear at distance exactly $\ell$ from the root, i.e., they lie on level $\ell$ of the tree. We now observe that
\[
\sum_{B\in\mathcal{B}_\ell}\textrm{charge}(C,B) \le \sum_{B\in\mathcal{B}_\ell}\frac{\eps}{h}\cdot \cost^\downarrow_B(C) \le \frac{\eps}{h}\cost(C),
\]
where the last inequality follows because if we have two bags $B,B'\in
\mathcal{B}_\ell$, then $E(G^\downarrow_B)\cap E(G^\downarrow_{B'})=\emptyset$:
note that every edge of $E(G^\downarrow_B)$ must be incident on a vertex $v$
that appears in a descendant of $B$, but not in $B$. By the properties of tree
decompositions, notably by the fact that $B$ is a separator of $G$, $v$ cannot
appear in $B'$ or any of its descendants.  Therefore none of its incident edges
are contained in $E(G^\downarrow_{B'})$.  Because
$\sum_{B\in\mathcal{B}_\ell}\cost^\downarrow_B(C)$ is the sum of costs of $C$
over disjoint sets of edges, the sum is a lower bound on the total cost of $C$.

To conclude, we observe that the total charge of $C$ is 
\[
\textrm{charge}(C)\le \sum_{\ell=0}^{h-1}\sum_{B\in\mathcal{B}_\ell}\textrm{charge}(C,B) \le \eps\cost(C).
\]
Therefore, summing over all components of $F^\star$, the total cost of the edges we have added according to \autoref{rule1} is at most $\eps\cdot\cost(F^\star)$.
\end{proof}

%
%

\subsubsection{Partitioning active terminals}

We are now ready to prove \autoref{lem:near-opt-1} for the near-optimal solution $F_\eps$ constructed above, for which we will compute the partitions $\zeta_B$ for all bags $B$.
We will use the following two claims for the active terminals~$A_B$ of the given bag $B$.

\begin{claim}\label{claim:1}%
If there exist $t_1,t_2\in A_B$ such that $\dist(t_1,t_2)\le \frac{\eps}{kh}\dist(t_1,B)$, then $t_1,t_2$ are in the same component of $F_\eps$.
\end{claim}

%
\begin{proof}%
Let $C_1$ be the component that contains $t_1$ in the optimal
solution $F^\star$ from which $F_\eps$ is constructed. We observe that
$\cost^\downarrow_{B}(C_1)\ge \dist(t_1,B)$, because~$C_1$ must contain a path
from $t_1$ to $B$ (as $t_1$ is active) and all the edges of this path are
contained in $E(G_{B}^\downarrow)$. If $t_2$ is contained in $C_2$ in $F^\star$
we therefore have, $\dist(C_1,C_2)\le \dist(t_1,t_2)\le
\frac{\eps}{kh}\dist(t_1,B) \le \frac{\eps}{kh} \cost^\downarrow_{B}(C_1)$.
Since $C_2$ must also intersect $B$, \autoref{rule1} implies that $C_1$ and
$C_2$ are contained in the same component of $F_\eps$.  \end{proof}   
%

\begin{claim}\label{claim:2} %
    Let $A\subseteq A_B$ and $d\ge 0$ be such that (i) there exists $b\in B$ such that for all $t\in A$ we have $\dist(t,B)=\dist(t,b)$ and $d\le \dist(t,B)\le 2d$ (ii) for all distinct $t,t'\in A$ we have $\dist(t,t')> \frac{\eps}{kh}d$ (iii)~$|A|\ge \frac{8k^2(k+1)h^2}{\eps^2}$. Then, there exists a component of $F_\eps$ that contains all terminals of $A$.
\end{claim}

%
\begin{proof}%
Consider an active component $C$ for $B$ of the optimum solution $F^\star$. We 
claim that $\cost^\downarrow_{B}(C)\ge |V(C)\cap A|\cdot \frac{\eps}{2kh}d$. 
%
%
%
    To see this, let $C^\downarrow_1,\ldots,C^\downarrow_\ell$ be the components of $C$ when restricting $C$ to $G^\downarrow_B$. Each component $C^\downarrow_i$ is a Steiner tree for the terminals in $V(C^\downarrow_i)\cap A$. Now consider a minimum spanning tree~$U$ on the metric closure derived from $G^\downarrow$ for this vertex set $V(C^\downarrow_i)\cap A$. It is well known that such a minimum spanning tree is a $2(1-\frac{1}{p})$-approximation of the optimum Steiner tree~\cite{KouMB81} on $p$ terminals, and thus $\cost(C^\downarrow_i)\geq\frac{1}{2(1-\frac{1}{p})}\cost(U)$ where $p=|V(C^\downarrow_i)\cap A|$. The distance between any two terminals of $V(C^\downarrow_i)\cap A$ in the given graph $G$ is more than $\frac{\eps}{kh}d$ by property~(ii) of the claim, and because the distance between such terminals can only be more in $G^\downarrow$, every edge of $U$ has cost more than $\frac{\eps}{kh}d$. This means that $\cost(U)\geq (p-1)\cdot\frac{\eps}{kh}d$, and we therefore get $\cost(C^\downarrow_i)\ge \frac{p-1}{2(1-\frac{1}{p})}\cdot \frac{\eps}{kh}d =\frac{p}{2}\cdot\frac{\eps}{kh}d$. Summing over all (vertex disjoint) components~$C^\downarrow_i$ we obtain the claimed inequality $\cost^\downarrow_{B}(C)\ge |V(C)\cap A|\cdot \frac{\eps}{2kh}d$.

    Because each terminal of $A$ belongs to an active component of $F^\star$, of which there are at most~$k+1$, there must exist an active component $C$ with $|V(C)\cap A|\geq\frac{|A|}{k+1}$, which by the above inequality and property~(iii) of the claim gives $\cost^\downarrow_{B}(C)\ge \frac{|A|}{k+1}\cdot \frac{\eps}{2kh}d\ge \frac{kh}{\eps}\cdot 4d$.
    %
    Now note that by property~(i), for all $t,t'\in A$ we have $\dist(t,t')\le 4d$, as we can use a path through $b$. So, if $C'$ is any of the other active components of $F^\star$ also containing a terminal of $A$, we have $\dist(C,C')\le 4d$. We therefore obtain $\dist(C,C')\le \frac{\eps}{kh}\cost^\downarrow_{B}(C)$, and according to \autoref{rule1}, $C$~and~$C'$ are part of the same component of $F_\eps$. In other words, all active terminals of $A$ are in components of $F^\star$ that lie in the one component of~$F_\eps$ containing~$C$.
\end{proof}
%

Intuitively, \autoref{claim:1} allows us to place terminals of $A$ which are very close to each other into the same set of the partition $\zeta_B$, as placing one terminal in a component forces the placement of the other. Thanks to this claim we can work with an appropriate net. 
If we find a large collection of such net points which also are roughly the same distance from the bag and closest to the same vertex of the bag, \autoref{claim:2} allows us to group them all together in the partition $\zeta_B$.
Armed with these tools, we can now prove the main lemma.

\begin{proof}[Proof of \autoref{lem:near-opt-1}]
To compute the partition $\zeta_B$ in polynomial time, we first partition the active terminals $A_B\cap B$ contained in the bag $B$. For this we simply add a set~$\{t\}$ for each $t\in A_B\cap B$ to $\zeta_B$, which adds at most $|B|\leq k+1$ sets to $\zeta_B$. Let now $A=A_B\setminus B$ be the remaining active terminals.

To partition $A$, let $d=\min_{t\in A_B\setminus B}\dist(t,B)$ and $D=\max_{t\in 
A_B\setminus B}\dist(t,B)$ be the minimum and maximum distances of these active 
terminals from the bag $B$. Then partition $A_B\setminus B$ into $|B|\le k+1$ 
sets $A_1, A_2,\ldots, A_{|B|}$, depending on the vertex of $B$ that is closest 
to each $t\in A$ (breaking ties arbitrarily). That is, for each $A_i$ there 
exists $b\in B$ such that for all $t\in A_i$ we have $\dist(t,B)=\dist(t,b)$. 
Consider now a set $A_i$ and further partition it into 
$r=\lceil\log_2\frac{D}{d}\rceil$ sets $A_{i,0}, A_{i,1},\ldots, A_{i,r-1}$, 
where $A_{i,j}$ contains all $t\in A_i$ such that $\dist(t,B)\in [2^jd, 
2^{j+1}d)$. Now (greedily) compute an $(\frac{\eps}{kh}2^jd)$-net $N_{i,j}$ of 
$A_{i,j}$. We observe that $N_{i,j}$ satisfies the first two conditions of 
\autoref{claim:2} for~$2^jd$, so if $|N_{i,j}|\ge \frac{8k^2(k+1)h^2}{\eps^2}$, 
then we add~$A_{i,j}$ as a set of our partition $\zeta_B$, remove the terminals 
of~$N_{i,j}$ from $A$ and continue the algorithm for the remaining terminals. 
Repeat the previous step for all~$i,j$ for which $N_{i,j}$ is sufficiently 
large. This contributes at most $(k+1)\lceil\log\frac{D}{d}\rceil$ sets to 
$\zeta_B$.

Suppose now that we are left with a set of terminals $A$ such that the procedure above fails to construct a sufficiently large net $N_{i,j}$ to apply \autoref{claim:2}. For every index pair $i,j$, each remaining terminal $t\in A_{i,j}$ is close enough to some net point $t'\in N_{i,j}$ such that we can apply \autoref{claim:1}. We therefore create a set in the partition~$\zeta_B$ for each $t'\in N_{i,j}$, placing into such a set those terminals of $A_{i,j}$ that are closest to~$t'$ (breaking ties arbitrarily). Since we cannot apply \autoref{claim:2} to the remaining sets $A_{i,j}$, each of the at most $(k+1)\lceil\log\frac{D}{d}\rceil$ nets $N_{i,j}$ has size less than $\frac{8k^2(k+1)h^2}{\eps^2}$, which implies $|\zeta_B|\leq O(\frac{k^4h^2}{\eps^2}\log\frac{D}{d})$.

Clearly the above procedure can be implemented in polynomial time, and the fact that every set of~$\zeta_B$ is contained in the same component of $F_\eps$ follows from \autoref{claim:1} and \autoref{claim:2}. Finally, any path in a graph with at most $n$ vertices has less than $n$ edges, so that $\frac{D}{d}<2n^2/\eps$, given that the ratio of the longest to the shortest edge is $2n/\eps$ (note that $d>0$ by definition). Hence the claimed bound of $|\zeta_B|\leq O(\frac{k^4h^2}{\eps^2}\log\frac{n}{\eps})$ follows.
\end{proof}

\begin{comment}
In order to get rid of the $\log(\frac{D}{d})$ factor in the bound of the above lemma,
we will use a recursive procedure based on the following definition...

\begin{definition}
Given $\eps>0$ and a bag $B$ of a tree decomposition of width $k$ and height $h$, 
let $X\subseteq B$ be a non-empty subset. We say that $X$ is \emph{well-separated} if 
$\diam(X)\leq (\frac{\eps}{kh})^2\cdot\dist(X,B\setminus X)$. We also declare the bag $B$
itself to be a well-separated set of $B$.
\end{definition}

Note that any singleton set of $B$ is well-separated.

\begin{lemma}
Given a bag $B$ of a tree decomposition, the family $\mathcal{X}_B$ of well-separated sets of $B$ is laminar.
\end{lemma}

The next lemma bounds the aspect ratio of the metric induced by the children of a 
well-separated set in the laminar family $\mathcal{X}_B$.

\begin{lemma}
Let $X\in\mathcal{X}_B$ be a well-separated set of a bag $B$.
If $d$ and $D$ are the respective minimum and maximum distance between any two children of 
$X$ in the laminar family $\mathcal{X}_B$, then $\frac{D}{d}\leq (\frac{kh}{\eps})^{2(k+1)}$.
\end{lemma}

In each recursion step we will be given a well-separated set $X$ and a subset $A$ of active 
terminals with an upper bound on their distances from $X$. 
If $X$ is not the bag itself and there is a terminal in $A$ that is far from $X$ then we 
will be able to partition $A$ directly, i.e., the recursion terminates.
Note that the following lemma can always be applied to a well-separated set that is a 
singleton (unless the bag itself is a singleton).

\begin{lemma}
Let $X\in\mathcal{X}_{B}\setminus\{B\}$ be a well-separated set of a bag $B$, and let
$A\subseteq A_i$ be a subset of the active terminals $A_i$ of $B$ such that
for all $t\in A$ we have $\dist(t,X)\leq\frac{\eps}{kh}\dist(X,B\setminus X)$.
If there exists a terminal $t\in A$ such that $\dist(t,X)\geq\frac{kh}{\eps}\diam(X)$,
then we can compute a partition $\xi$ of $A$ in polynomial time, such that every set of $\xi$
belongs to the same component of $F_\eps$ and~$|\xi|\leq\cdot(\frac{kh}{\eps})^{O(1)}$.
\end{lemma}

Finally we handle the case when there is no terminal far from a well-separated set.

\begin{lemma}
Let $X\in\mathcal{X}_{B}\setminus\{B\}$ be a well-separated set of a bag $B$, and let
$A\subseteq A_i$ be a subset of the active terminals $A_i$ of $B$ such that
for all $t\in A$ we have $\dist(t,X)\leq\frac{kh}{\eps}\diam(X)$.
Then we can recursively compute a partition $\xi$ of $A$ in polynomial time, such that every
set of $\xi$ belongs to the same component of $F_\eps$ 
and~$|\xi|\leq\cdot(\frac{kh}{\eps})^{O(1)}$.
\end{lemma}

We use the above recursive procedure to compute the partition $\zeta_i$ of a bag $B$
as given by \autoref{lem:near-opt-1}.
\end{comment}

